\section{附錄} % Main chapter title
 %\section{人物}

\subsection{四福音大纲}
\subsubsection{马太福音}

\subsubsection{马可福音}


\subsubsection{路加福音}
\begin{table}[htbp]
  \centering
  \caption{耶稣的事工与生平概要}
  \begin{tabular}{cll}
    \toprule
    \textbf{编号} & \textbf{事工与主题描述} & \textbf{章节范围} \\
    \midrule
    % 原 \subsubsection{耶稣的服事 1-19:26} 转化为表格的分类大项
    \multicolumn{3}{l}{\textbf{一、耶稣的服事}} \\
    \quad 1.1 & 耶稣服事前 & 1--4:15 \\
    \quad 1.2 & 耶稣呼召门徒 & 4:16--6:11 \\
    \quad 1.3 & 耶稣拣选十二门徒 & 6:12--8 \\
    \quad 1.4 & 耶稣差遣十二门徒 & 9 \\
    \quad 1.5 & 耶稣差遣七十二门徒 & 10:12--17:10 \\
    \quad 1.6 & 耶稣往耶路撒冷途中的服事 & 17:11--19:26 \\
    \midrule
    % 原 \subsubsection{耶稣的受死和复活 19:27-24:53}
    \multicolumn{3}{l}{\textbf{二、耶稣的受死和复活}} \\
    \quad 1.7 & 耶稣的受死和复活 & 19:27--24:53 \\
    \bottomrule
  \end{tabular}
\end{table}



\subsubsection{約翰福音}
\begin{table}[htbp]
  \centering
  \caption{耶稣在耶路撒冷的事工与受死复活概要}
  \begin{tabular}{cll}
    \toprule
    \textbf{编号} & \textbf{事工与主题描述} & \textbf{章节范围} \\
    \midrule
    % 大类一
    \multicolumn{3}{l}{\textbf{一、前四次在耶路撒冷————耶稣的服事 (1--11)}} \\
    \quad 1.1 & 逾越節第一次到耶路撒冷 & 1--3:21 \\
    \quad 1.2 & 節期第二次到耶路撒冷 & 3:22--5 \\
    \quad 1.3 & 住棚節第三次到耶路撒冷 & 6--10:21 \\
    \quad 1.4 & 修殿節第四次到耶路撒冷 & 10:22--11 \\
    \midrule
    % 大类二
    \multicolumn{3}{l}{\textbf{二、第五次到耶路撒冷————耶稣的受死和复活 (12--21)}} \\
    \quad 2.1 & 对犹太人的最后讲话 & 13 \\
    \quad 2.2 & 对门徒的讲话 & 14--17 \\
    \quad 2.3 & 耶稣的被捕受审受死 & 18--19:16 \\
    \quad 2.4 & 耶稣的复活 & 19:17--21 \\
    \bottomrule
  \end{tabular}
\end{table}





\subsection{安息日所行的八件事}
\begin{table}[H]
\centering
\begin{tabular}{p{4.37cm}|p{1.5cm}|p{1.5cm}|p{1.5cm}|p{1.350cm}|p{1.0cm}|p{3.7cm}}\hline\hline
\multicolumn{1}{c|}{事件}&\multicolumn{1}{c|}{太}&\multicolumn{1}{c|}{可}&\multicolumn{1}{c|}{路}&\multicolumn{1}{c|}{约}&\multicolumn{1}{c|}{神迹}&\multicolumn{1}{c}{被责备}\\\hline
迦百农\textcolor{red}{会堂}赶逐污鬼&&1:23-26&4:31-37&&是&\\\hline
迦百农治好彼得岳母&8:14-17&1:29-31&4:38-41& &是&\\\hline
门徒在麦地掐麦穗&&2:23-28&6:1-5& &不是&法利赛\\\hline
在\textcolor{red}{会堂}治好萎缩的右手&12:9-14&3:1-6&6:6-11& &是&文士，法利赛，希律党\\\hline
在\textcolor{red}{会堂}医治驼背的女人&&&13:10-17& &是&管会堂的人\\\hline
法利赛\textcolor{red}{首领家}治水臌病人&&&14:1-6& &是&律法师，法利赛人\\\hline
耶路撒冷毕士大池治瘫子&&& &5:1-18 &是&犹太人\\\hline
耶路撒冷治生来瞎眼的&&&&9&是&法利赛人，犹太人\\\hline
\end{tabular}
\end{table}

\subsection{耶穌和人不同的看法}
\begin{table}[htbp]
\centering
\caption{对比主题汇总（网格版）}
\begin{tabular}{c|l|l}\hline
\textbf{编号} & \textbf{内容} & \textbf{章节范围} \\\hline
1.1 & 去叫醒拉撒路 vs 和他同被石头打死 & 约 11:1-16 \\\hline
1.2 & 耶稣说的复活 vs 马大理解的复活 & 约 11:21-40 \\\hline
1.3 & 拉撒路复活信耶稣 vs 告诉法利赛人 & 约 11:45-46 \\\hline
1.4 & 膏抹主 vs 挂念穷人 & 约 12:4-8 \\\hline
1.5 & 鬼服了你们欢喜 vs 名记录在天上欢喜 & 路 10:17--20 \\\hline
1.6 & 在耶稣脚前听道 vs 为许多的事思虑烦扰 & 路 10:38--42 \\\hline
1.7 & 死而复活、失而又得的快乐 vs 大儿子不肯进去生气 & 路 15:11-32 \\\hline
1.8 & 税吏的祷告 vs 自高的法利赛人 & 路 18:9--14 \\\hline
1.9 & 亚伯拉罕的子孙 vs 罪人税吏撒该 & 路 19:1--10 \\\hline
\end{tabular}
\end{table}





\subsection{耶穌问的问题}
\begin{table}[H]
  \centering  
    \begin{tabular}{p{4.72em}|p{40.72em}} \toprule
 经文  &\textbf{问题} \\ \midrule
太9:1-8 &耶穌上了船，渡過海，來到自己的城裡。 2 有人用褥子抬著一個癱子到耶穌跟前來。耶穌見他們的信心，就對癱子說：「小子，放心吧，你的罪赦了！」 3 有幾個文士心裡說：「這個人說僭妄的話了！」 4 耶穌知道他們的心意，就說：「\textcolor{red}{你們為什麼心裡懷著惡念呢？ 5 或說『你的罪赦了』，或說『你起來行走』，哪一樣容易呢？} 6 但要叫你們知道，人子在地上有赦罪的權柄。」就對癱子說：「起來，拿你的褥子回家去吧！」 7 那人就起來，回家去了。 8 眾人看見都驚奇，就歸榮耀於神，因為他將這樣的權柄賜給人。   \\ \midrule 
太9:27-31&耶穌從那裡往前走，有兩個瞎子跟著他，喊叫說：「大衛的子孫，可憐我們吧！」 28 耶穌進了房子，瞎子就來到他跟前，耶穌說：「\textcolor{red}{你們信我能做這事嗎？}」他們說：「主啊，我們信。」 29 耶穌就摸他們的眼睛，說：「照著你們的信給你們成全了吧！」 30 他們的眼睛就開了。耶穌切切地囑咐他們，說：「你們要小心，不可叫人知道。」 31 他們出去，竟把他的名聲傳遍了那地方。\\\midrule

太16:13-20&耶穌到了愷撒利亞-腓立比的境內，就問門徒說：「\textcolor{red}{人說我人子是誰？}」 14 他們說：「有人說是施洗的約翰，有人說是以利亞，又有人說是耶利米或是先知裡的一位。」 15 耶穌說：「\textcolor{red}{你們說我是誰？}」 16 西門彼得回答說：「你是基督，是永生神的兒子。」 17 耶穌對他說：「西門‧巴約拿，你是有福的，因為這不是屬血肉的指示你的，乃是我在天上的父指示的。 18 我還告訴你：你是彼得，我要把我的教會建造在這磐石上，陰間的權柄 b不能勝過他。 19 我要把天國的鑰匙給你，凡你在地上所捆綁的，在天上也要捆綁；凡你在地上所釋放的，在天上也要釋放。」 20 當下，耶穌囑咐門徒，不可對人說他是基督。\\\midrule

%太17:1&\\\midrule

    \bottomrule
    \end{tabular}%
 %   \caption{被耶穌誇獎的人}
  \label{tab:addlabel}%
\end{table}%

\subsection{问的耶穌问题}
\begin{table}[H]
  \centering  
    \begin{tabular}{p{4.72em}|p{40.72em}} \toprule
 经文& \textbf{问题}  \\ \midrule
太9:1-8&耶穌從那裡往前走，看見一個人，名叫馬太，坐在稅關上，就對他說：「你跟從我來！」他就起來跟從了耶穌。耶穌在屋裡坐席的時候，有好些稅吏和罪人來，與耶穌和他的門徒一同坐席。 11 法利賽人看見，就對耶穌的門徒說：「\textcolor{red}{你們的先生為什麼和稅吏並罪人一同吃飯呢？}」 12 耶穌聽見就說：「康健的人用不著醫生，有病的人才用得著。 13 經上說：『我喜愛憐恤，不喜愛祭祀』，這句話的意思你們且去揣摩；我來本不是召義人，乃是召罪人。」 \\ \midrule  
太9:14-18& 那時，約翰的門徒來見耶穌，說：「\textcolor{red}{我們和法利賽人常常禁食，你的門徒倒不禁食，這是為什麼呢？}」 15 耶穌對他們說：「新郎和陪伴之人同在的時候，陪伴之人豈能哀慟呢？但日子將到，新郎要離開他們，那時候他們就要禁食。 16 沒有人把新布補在舊衣服上，因為所補上的反帶壞了那衣服，破的就更大了。 17 也沒有人把新酒裝在舊皮袋裡，若是這樣，皮袋就裂開，酒漏出來，連皮袋也壞了。唯獨把新酒裝在新皮袋裡，兩樣就都保全了。」   \\\midrule
太11:2-19&約翰在監裡聽見基督所做的事，就打發兩個門徒去，問他說：「\textcolor{red}{那將要來的是你嗎？還是我們等候別人呢？}」耶穌回答說：「你們去，把所聽見、所看見的事告訴約翰，就是：瞎子看見，瘸子行走，長大痲瘋的潔淨，聾子聽見，死人復活，窮人有福音傳給他們。凡不因我跌倒的，就有福了！」他們走的時候，耶穌就對眾人講論約翰說：「你們從前出到曠野是要看什麼呢？要看風吹動的蘆葦嗎？你們出去到底是要看什麼？要看穿細軟衣服的人嗎？那穿細軟衣服的人是在王宮裡。你們出去究竟是為什麼？是要看先知嗎？我告訴你們，是的，他比先知大多了。經上記著說：『我要差遣我的使者在你前面預備道路』，所說的就是這個人。我實在告訴你們：凡婦人所生的，沒有一個興起來大過施洗約翰的；然而天國裡最小的比他還大。從施洗約翰的時候到如今，天國是努力進入的，努力的人就得著了。因為眾先知和律法說預言到約翰為止；你們若肯領受，這人就是那應當來的以利亞。有耳可聽的，就應當聽！我可用什麼比這世代呢？好像孩童坐在街市上招呼同伴說：『我們向你們吹笛，你們不跳舞！我們向你們舉哀，你們不捶胸！』約翰來了，也不吃也不喝，人就說他是被鬼附著的。人子來了，也吃也喝，人又說他是貪食好酒的人，是稅吏和罪人的朋友。但智慧之子總以智慧為是。\\\midrule
太12:9-21&耶穌離開那地方，進了一個會堂。 10 那裡有一個人枯乾了一隻手。有人問耶穌說：「\textcolor{red}{安息日治病可以不可以？}」意思是要控告他。 11 耶穌說：「你們中間誰有一隻羊當安息日掉在坑裡，不把牠抓住拉上來呢？12 人比羊何等貴重呢！所以，在安息日做善事是可以的。」 13 於是對那人說：「伸出手來！」他把手一伸，手就復了原，和那隻手一樣。\\\midrule
太13:10-16&門徒進前來問耶穌說：「\textcolor{red}{對眾人講話為什麼用比喻呢？}」 耶穌回答說：「因為天國的奧祕只叫你們知道，不叫他們知道。凡有的，還要加給他，叫他有餘；凡沒有的，連他所有的也要奪去。所以我用比喻對他們講，是因他們看也看不見，聽也聽不見，也不明白。在他們身上，正應了以賽亞的預言說：『你們聽是要聽見，卻不明白；看是要看見，卻不曉得。 15 因為這百姓油蒙了心，耳朵發沉，眼睛閉著，恐怕眼睛看見，耳朵聽見，心裡明白，回轉過來，我就醫治他們。』 16 但你們的眼睛是有福的，因為看見了；你們的耳朵也是有福的，因為聽見了。\\\midrule
太15:1-9&那時，有法利賽人和文士從耶路撒冷來見耶穌，說： 2 「\textcolor{red}{你的門徒為什麼犯古人的遺傳呢？}因為吃飯的時候，他們不洗手。」 3 耶穌回答說：「你們為什麼因著你們的遺傳，犯神的誡命呢？ 4 神說：『當孝敬父母』，又說：『咒罵父母的，必治死他。』 5 你們倒說，無論何人對父母說：『我所當奉給你的已經做了供獻』， 6 他就可以不孝敬父母。這就是你們藉著遺傳，廢了神的誡命！ 7 假冒為善的人哪！以賽亞指著你們說的預言是不錯的，他說： 8 『這百姓用嘴唇尊敬我，心卻遠離我。 9 他們將人的吩咐當做道理教導人，所以拜我也是枉然。』」\\\midrule
太17:9-13&下山的時候，耶穌吩咐他們說：「人子還沒有從死裡復活，你們不要將所看見的告訴人。」 10 門徒問耶穌說：「\textcolor{red}{文士為什麼說以利亞必須先來？}」 11 耶穌回答說：「以利亞固然先來，並要復興萬事。 12 只是我告訴你們：以利亞已經來了，人卻不認識他，竟任意待他。人子也將要這樣受他們的害。」 13 門徒這才明白耶穌所說的是指著施洗的約翰。\\\midrule
太17:19-21&門徒暗暗地到耶穌跟前說：「\textcolor{red}{我們為什麼不能趕出那鬼呢？}」 20 耶穌說：「是因你們的信心小。我實在告訴你們：你們若有信心像一粒芥菜種，就是對這座山說：『你從這邊挪到那邊！』，它也必挪去；並且你們沒有一件不能做的事了。 21 至於這一類的鬼，若不禱告、禁食，他就不出來 a。」\\\midrule

    \bottomrule
    \end{tabular}%
  %  \caption{被耶穌誇獎的人}
  \label{tab:addlabel}%
\end{table}%


\subsection{法利賽人/撒都該人的错误}
\begin{table}[H]
  \centering  
    \begin{tabular}{p{4.72em}|p{40.72em}} \toprule
 经文& \textbf{问题}  \\ \midrule
 太15:1-9/路7:1-16俗手吃饭和出口的汙穢&那時，有法利賽人和文士從耶路撒冷來見耶穌，說： 2 「你的門徒為什麼犯古人的遺傳呢？因為吃飯的時候，他們不洗手。」 3 耶穌回答說：「\textcolor{red}{你們為什麼因著你們的遺傳，犯神的誡命呢？ 4 神說：『當孝敬父母』，又說：『咒罵父母的，必治死他。』 5 你們倒說，無論何人對父母說：『我所當奉給你的已經做了供獻』， 6 他就可以不孝敬父母。這就是你們藉著遺傳，廢了神的誡命！} 7 假冒為善的人哪！以賽亞指著你們說的預言是不錯的，他說： 8 『這百姓用嘴唇尊敬我，心卻遠離我。 9 他們將人的吩咐當做道理教導人，所以拜我也是枉然。』」\\\midrule
責法利賽撒都該人求天上神迹 穢&\\ \midrule
防法利賽/撒都該/希律的酵 &\\ \midrule
    \bottomrule
    \end{tabular}%
  %  \caption{法利賽人的错误教导}
  \label{tab:addlabel}%
\end{table}%


\subsection{被耶穌誇獎的人}
\begin{table}[H]
  \centering  
    \begin{tabular}{l|r|r|r|l} \toprule
  \textbf{讨主喜悦的人} & \multicolumn{1}{l|}{\textbf{太}} & \multicolumn{1}{l|}{\textbf{可}} & \multicolumn{1}{l|}{\textbf{路}} & \textbf{约} \\ \midrule
真以色列人拿但業 &       &       &  &1:45-51  \\ \midrule 
施洗约翰  &       &       &       &  \\\midrule
四個抬摊子的人，使癱子得医治 & \multicolumn{1}{l|}{9:1-8} & \multicolumn{1}{l|}{2:1-12} & \multicolumn{1}{l|}{5:17-26} &  \\ \midrule
百夫长的信心使仆人得医治 & \multicolumn{1}{l|}{8:5-13} &       & \multicolumn{1}{l|}{7:1-10} &  \\ \midrule
十個大麻風中的撒瑪利亞人 & &  & \multicolumn{1}{l|}{17:11-19} &\\\midrule
瞎子巴低买 &       & \multicolumn{1}{l|}{10:46-52} & \multicolumn{1}{l|}{18:35-43} &  \\ \midrule
亞伯拉罕的子孫稅吏撒該&       & \multicolumn{1}{l|}{10:46-52} & \multicolumn{1}{l|}{15:1-10} &  \\ \midrule
十字架上的强盗&       & \multicolumn{1}{l|}{10:46-52} & \multicolumn{1}{l|}{15:1-10} &  \\ \midrule

妓女用眼泪给主洗脚 &       &       & \multicolumn{1}{l|}{} &  \\ \midrule
血露妇人因信心得醫治&\multicolumn{1}{l|}{9:20-22}&\multicolumn{1}{l|}{5:24-34}&\multicolumn{1}{l|}{8:43-48}&\\\midrule
叙利非尼基妇人女儿得医治 & \multicolumn{1}{l|}{15:21-28} & \multicolumn{1}{l|}{6:24-30} &       &  \\ \midrule
穷寡妇的奉献 &       & \multicolumn{1}{l|}{12:41-44} & \multicolumn{1}{l|}{21:1-4} &  \\  \midrule
玛丽亚得上好的福分 & \multicolumn{1}{l|}{} & \multicolumn{1}{l|}{ } & \multicolumn{1}{l|}{10:38-42 } &  \\ \midrule
伯大尼的玛丽亚膏主 & \multicolumn{1}{l|}{26:6-13} & \multicolumn{1}{l|}{14:3-9} & \multicolumn{1}{l|}{ } & 12:1-8 \\
 %   \midrule
 %  約瑟和尼哥底母埋葬耶穌 & \multicolumn{1}{l|}{27：57-60}  & \multicolumn{1}{l|}{15：42-46} & \multicolumn{1}{l|}{23：50-53} & 19：38-42 \\
  %  \midrule
   % 抹大拉的玛利亚在坟墓哭主 &       &       &       & 20：11-18 \\
    \bottomrule
    \end{tabular}%
    \caption{被耶穌誇獎的人}
  \label{tab:addlabel}%
\end{table}%


\subsection{四位馬利亞}
\begin{table}[H]
  \centering
    \begin{tabular}{p{10.22em}|p{8.72em}|p{6.055em}|p{7.61em}|p{5.055em}|p{4.72em}}    \toprule
    \multicolumn{1}{p{10.22em}|}{人物} & 馬太福音  & 馬可福音  & 路加福音  & 約翰福音  & \multicolumn{1}{l}{其他書卷} \\    \midrule
    \multicolumn{1}{p{10.22em}|}{拿撒勒的馬利亞（主耶穌的母親）} & 1:16,18-25;2:10,11; 12:46-53;13:53-58 & 3:31-35;6:1-6 &1:26-56;2:1-52; 8:19-21;11:27-28 &2:1-12; 19:25-27&徒1:14，加4:4 \\    \midrule
    \multicolumn{1}{p{10.22em}|}{耶路撒冷的馬利亞(革羅罷妻子/馬可母親)} & 27:55-56,61;28:1-10 & 15:40-41; 16:1-8,12-13 & 23:49,55-56;24:1-11;22-24 & 19:25 & 西4:10，彼前5:13 \\    \midrule
    抹大拉的馬利亞 & 27:55-56,61;28:1-10 & 15:40-41,47; 16:1-11 & 8:2-3,23:49,55-56,24:1-11 & 19:25;20:1-18 &  \\    \midrule
    伯大尼的馬利亞 & 26:6-13  & 14:3-9 &10:38-42 &11:1-20,28-45;12:1-8 &  \\
    \bottomrule
    \end{tabular}%
     \caption{四位馬利亞}
  \label{tab:addlabel}%
\end{table}%


%\section{對比}

\subsection{法利賽人求神蹟}

\begin{table}[H]
  \centering
  \begin{tabular}{p{4.2em}|p{22.3em}|p{15.3em}}  
  \toprule
    &   太12:38-42  
    &   太16:1-4   \\    \midrule
   人群 & 有幾個文士和法利賽人& 法利賽人和撒都該人來試探耶穌  \\    \midrule
    問題 & 夫子，我們願意你顯個神蹟給我們看。 & 請他從天上顯個神蹟給他們看。 \\
    \midrule
   主的回答 & 一個邪惡淫亂的世代求看神蹟，除了先知約拿的神蹟以外，再沒有神蹟給他們看。約拿三日三夜在大魚肚腹中，人子也要這樣三日三夜在地裡頭。當審判的時候，尼尼微人要起來定這世代的罪，因為尼尼微人聽了約拿所傳的就悔改了。看哪，在這裡有一人比約拿更大！當審判的時候，南方的女王要起來定這世代的罪，因為他從地極而來，要聽所羅門的智慧話。看哪！在這裡有一人比所羅門更大。& 晚上天發紅，你們就說：『天必要晴。』早晨天發紅，又發黑，你們就說：『今日必有風雨。」你們知道分辨天上的氣色，倒不能分辨這時候的神蹟。一個邪惡淫亂的世代求神蹟，除了約拿的神蹟以外，再沒有神蹟給他看。」耶穌就離開他們去了。\\    \bottomrule
    \end{tabular}%
  \label{tab:addlabel}%
 \caption{法利賽人求神蹟}  
\end{table}%


\subsection{哀哭切齿的人}
\begin{table}[H]
  \centering
  \begin{tabular}{p{4.2em}|p{44.3em}}  \toprule
 经文 &   内容   \\ \midrule
伯16:9,10&主發怒撕裂我,逼迫我,向我切齒;我的敵人怒目看我。他們向我開口,打我的臉羞辱我,聚會攻擊我。 \\ \midrule
詩35:16&他們如同席上好戲笑的狂妄人，向我咬牙。 \\ \midrule
詩37:15-16&我在患難中，他們卻歡喜，大家聚集。我所不認識的那些下流人聚集攻擊我，他們不住地把我撕裂。惡人設謀害義人，又向他咬牙。 \\ \midrule
詩112:9-10&他施捨錢財，賙濟貧窮，他的仁義存到永遠。他的角必被高舉，大有榮耀。
10 惡人看見便惱恨，必咬牙而消化，惡人的心願要歸滅絕。\\ \midrule
哀2:16&你(錫安城)的仇敵都向你大大張口，他們嗤笑又切齒，說：「我們吞滅她！這真是我們所盼望的日子臨到了，我們親眼看見了！」 \\ \midrule
太8:11-12 &我又告诉你们:从东从西，将有许多人来，在天国里与亚伯拉罕、以撒、雅各一同坐席；唯有本国的子民，竟被赶到\textcolor{red}{外边黑暗里}去，在那里必要哀哭切齿了。\\    \midrule
太13:40-42 & 将稗子薅出来用火焚烧，世界的末了也要如此。人子要差遣使者，把一切叫人跌倒的和作恶的，从他国里挑出来，丢在\textcolor{red}{火炉里}，在那里必要哀哭切齿了。\\ \midrule   
太13:49，50 & 天国又好像网撒在海里，聚拢各样水族。网既满了，人就拉上岸来，坐下，拣好的收在器具里，将不好的丢弃了。世界的末了也要这样。天使要出来，从义人中把恶人分别出来，丢在\textcolor{red}{火炉里}，在那里必要哀哭切齿了。\\ \midrule

太22:10-14 &天国好比一个王为他儿子摆设娶亲的筵席，就打发仆人去，请那些被召的人来赴席，他们却不肯来。王又打发别的仆人，说：‘你们告诉那被召的人：我的筵席已经预备好了，牛和肥畜已经宰了，各样都齐备，请你们来赴席。那些人不理就走了：一个到自己田里去，一个做买卖去，其余的拿住仆人，凌辱他们，把他们杀了。王就大怒，发兵除灭那些凶手，烧毁他们的城。于是对仆人说：‘喜筵已经齐备，只是所召的人不配。所以你们要往岔路口上去，凡遇见的，都召来赴席。’ 那些仆人就出去，到大路上，凡遇见的，不论善恶都召聚了来，筵席上就坐满了客。王进来观看宾客，见那里有一个没有穿礼服的，就对他说：‘朋友，你到这里来怎么不穿礼服呢？’那人无言可答。于是王对使唤的人说：‘捆起他的手脚来，把他丢在\textcolor{red}{外边的黑暗里}，在那里必要哀哭切齿了。’因为被召的人多，选上的人少。”\\ \midrule

太24:45-51 & 谁是忠心有见识的仆人，为主人所派管理家里的人，按时分粮给他们呢？主人来到，看见他这样行，那仆人就有福了。我实在告诉你们：主人要派他管理一切所有的。倘若那恶仆心里说‘我的主人必来得迟’，就动手打他的同伴，又和酒醉的人一同吃喝，在想不到的日子，不知道的时辰，那仆人的主人要来，重重地处治他，定他\textcolor{red}{和假冒为善的人}同罪，在那里必要哀哭切齿了。\\ \midrule

太25:28-30 &夺过他这一千来，给那有一万的！因为凡有的，还要加给他，叫他有余；没有的，连他所有的也要夺过来。把这无用的仆人丢在\textcolor{red}{外面黑暗里}，在那里必要哀哭切齿了。\\  \midrule 
路13:22-29&耶穌往耶路撒冷去，在所經過的各城各鄉教訓人。有一個人問他說：「主啊，得救的人少嗎？」耶穌對眾人說：「你們要努力進窄門。我告訴你們：將來有許多人想要進去，卻是不能。及至家主起來關了門，你們站在外面叩門，說：『主啊，給我們開門！』他就回答說：『我不認識你們，不曉得你們是哪裡來的。』那時，你們要說：『我們在你面前吃過、喝過，你也在我們的街上教訓過人。』他要說：『我告訴你們，我不曉得你們是哪裡來的。你們這一切作惡的人，離開我去吧！』你們要看見亞伯拉罕、以撒、雅各和眾先知都在神的國裡，你們卻\textcolor{red}{被趕到外面}，在那裡必要哀哭切齒了。從東、從西、從南、從北將有人來，在神的國裡坐席。只是有在後的將要在前，有在前的將要在後。」\\    \bottomrule
    \end{tabular}%
  \label{tab:addlabel}%
 \caption{哀哭切齿}  
\end{table}%



\subsection{預言的應驗}
\subsubsection{耶穌說的預言}




\begin{table}[htbp]
\centering
\caption{耶稣教导与预言主题汇总}
\small % Keeps the text size readable but compact for A4
\begin{tabularx}{\textwidth}{|>{\centering\arraybackslash}X|>{\arraybackslash}X|>{\raggedright\arraybackslash}X|}\hline
\textbf{主题} & \textbf{内容} & \textbf{章节范围} \\\hline

咒詛無花果樹 & 咒詛無花果樹 & 太21:18--22；可11:12--14；可11:20--26；可13:28--37；路13:6--9 \\\hline

\multirow{5}{*}{葡萄園} 
& 葡萄園工人 & 太20:1--16 \\\cline{2-3}
& 兩個兒子 & 太21:28--32 \\\cline{2-3}
& 凶惡園戶 & 太21:33--46；可12:1--12 \\\cline{2-3}
& 葡萄園中的無花果樹 & 路13:6--9 \\\hline

你所賜給我的人，我沒有失落一個 & 你所賜給我的人，我沒有失落一個 & 約18:9 \\\hline

拆毀聖殿三日建造／約拿三日三夜在大魚肚腹 & 拆毀聖殿三日建造／約拿三日三夜在大魚肚腹 & 約2:19--22；太12:40；可15:29 \\\hline

預言彼得三次不認主 & 預言彼得三次不認主 & 太26:34；可14:30；路22:34；約13:38 \\\hline

預言猶大出賣耶穌 & 預言猶大出賣耶穌 & 太26:21--25；可14:18--21；路22:21--23；約13:21--30 \\\hline

吩咐門徒預備驢 & 吩咐門徒預備驢 & 太21:1--7；可11:1--7；路19:29--35 \\\hline

吩咐門徒預備逾越節筵席 & 吩咐門徒預備逾越節筵席 & 太26:17--19；可14:12--16；路22:7--13 \\\hline

預言有人在沒死前見主榮神國 & 預言有人在沒死前見主榮神國 & 太16:28；可9:1；路9:27 \\\hline
\end{tabularx}
\end{table}





\subsubsection{聖經預言的應驗}

\begin{table}[htbp]
\centering
\caption{耶稣生平、受难与旧约预言应验汇总}
\small % Keeps text compact to fit on A4 margins safely
\begin{tabular}{|p{3.3cm}|p{4.5cm}|p{4.6cm}|}\hline
\textbf{主题} & \textbf{内容 / 应验经文} & \textbf{相关章节与经文背景} \\\hline

骑着小驴驹-和平的君王 & 耶稣得了一个驴驹，就骑上。 & 约12:14--15；赛62:11；亚9:9 \\\hline

你从婴孩和喫奶的口中完成了赞美的话 & 耶稣在圣殿中回应祭司长的责备。 & 太21:12--16；诗8:2 \\\hline

我们所传的有谁信呢 & 虽看见神迹，他们仍是不信。 & 约12:37--41；赛53:1；赛6:10 \\\hline

基督是大卫的子孙 & 论到基督的身份与大卫的后裔。 & 预言主题（太22:41--46） \\\hline

匠人所弃的石头已成了房角的头块石头 & 耶稣论到园户的比喻与房角石。 & 太21:42；可12:10；诗118:22 \\\hline

血田 & 犹大出卖耶稣的银钱用以买田。 & 亚11:12--13；太27:3--10 \\\hline

正月十日逾越节羊羔预备好 & 预备逾越节，耶稣荣耀进入耶路撒冷。 & 预表/进入耶路撒冷 \\\hline

正月十四逾越节羊羔-被杀 & 耶稣作为逾越节的羊羔被钉十字架。 & 预表/受难 \\\hline

初熟节首生-复活 & 耶稣从死里复活，成为初熟的果子。 & 预表/复活（林前15:20） \\\hline

击打牧人，羊就分散 & 预言门徒在祂受难当夜都会跌倒。 & 太26:31；可14:27；亚13:7 \\\hline

被列在罪犯之中 & 耶稣叮嘱门徒预备，受难时被列在罪犯中。 & 路22:35--38；赛53:12；可15:28 \\\hline

祂被欺压在受苦的时候却不开口 & 祂像羊羔被牵到宰杀之地，默默无声。 & 赛53:7 \\\hline

祂被挂在木头上，亲身担当了我们的罪 & 耶稣被钉十字架，代替世人的罪孽。 & 预言主题（彼前2:24） \\\hline

同我吃饭的人用脚踢我 & 预言同台吃饭的人（犹大）会出卖祂。 & 约13:18--20；诗41:9 \\\hline

因祂受的鞭伤我们得医治 & 耶稣在受审时被罗马兵丁残忍鞭打。 & 预言主题/被鞭打（赛53:5） \\\hline

分了我的外衣，为我的里衣拈鬮 & 兵丁在十字架下分了耶稣的外衣和里衣。 & 约19:23--24；诗22:18 \\\hline

骨头一根也不可折断 & 兵丁打断旁人的腿，但见耶稣已死便没有折断祂的骨头。 & 约19:33--36；出12:46；诗34:20 \\\hline

仰望自己所扎的人 & 兵丁用枪扎耶稣的肋旁，有血和水流出来。 & 约19:34--37；亚12:10 \\\hline

我渴了 & 耶稣为了使经上的话应验，在十架上说。 & 约19:28--30；诗69:21 \\\hline

无瑕疵的羊羔 & 耶稣没有做过一件不好的事。 & 十架强盗见证（路23:41）；出12:5 \\\hline

与财主同葬 & 祂虽然未行强暴，死的时候却与财主同葬。 & 赛53:9；太27:57--60 \\\hline

\end{tabular}
\end{table}










\subsubsection{舊約聖經預言表格}
\begin{table}[htbp]
  \centering
    \begin{tabular}{|p{10.72em}|p{9.11em}|p{4.445em}|l|l|p{8.39em}|}
    \toprule
    \multicolumn{1}{|l|}{\textbf{事件}} & \multicolumn{1}{l|}{\textbf{內容}} & \multicolumn{1}{l|}{\textbf{馬太福音}} & \multicolumn{1}{l|}{\textbf{馬可福音}} & \multicolumn{1}{l|}{\textbf{路加福音}} &   \multicolumn{1}{l|}{\textbf{舊約聖經}} \\
    \midrule
               施洗約翰誕生 & 先鋒要先來 & 11:10 & \multicolumn{1}{p{4.39em}|}{1:2} & \multicolumn{1}{p{4.165em}|}{7:27} &瑪3:1 \\
    \midrule
    施洗約翰的工作 & 曠野之聲  & 3:3  & \multicolumn{1}{p{4.39em}|}{1:3} & \multicolumn{1}{p{4.165em}|}{ 3:4} &        賽40:3-5 \\
    \midrule
    施洗約翰第二次的見證 & 贖罪祭   &   &       &    1:29   &        賽53:10 \\
  \midrule      
    \multirow{3}{*}{天使向馬利亞顯現} & 女人後裔  & 1:18 &       &       &        創3:15 \\
\cmidrule{2-6}    \multicolumn{1}{|r|}{} & 童女所生  & 1:22-23 &       &       &        賽7:14 \\
\cmidrule{2-6}    \multicolumn{1}{|r|}{} & 稱為以馬內利 & 1:22-23 &       &       &        賽7:14 \\
    \midrule
     \multirow{2}{*}{耶穌基督降生}  & 生在伯利恆 & 2:1，6 &       & \multicolumn{1}{p{4.165em}|}{2：4-7} &        彌5:2-3 \\
   % \midrule
   \cmidrule{2-6}    \multicolumn{1}{|r|}{}    & 希律屠嬰   & 2:16-18 &       &  &        耶利米31：15 \\
    \midrule
    \multirow{4}{*}{耶穌被列入家譜} & 亞伯拉罕的後裔 & 1:1-2 &       & \multicolumn{1}{p{4.165em}|}{3:34} &        創22:18，26:4 \\
\cmidrule{2-6}    \multicolumn{1}{|r|}{} & 猶大支派  &1:2-3 &       & \multicolumn{1}{p{4.165em}|}{3:33-34} &        創49:8-10 \\
\cmidrule{2-6}    \multicolumn{1}{|r|}{} & 耶西的根  & 1:5-6 &       & \multicolumn{1}{p{4.165em}|}{3:32} &        賽11:1，11:10 \\
\cmidrule{2-6}    \multicolumn{1}{|r|}{} & 大衛子孫  & 1:1，10 &       & \multicolumn{1}{p{4.165em}|}{3:31-32} &        撒下7:12-16，代上17:11-14，詩89:3-4，132：11，耶23：5 \\
    \midrule
    約瑟帶耶穌等逃往埃及 & 從埃及召出 & 2:15 &       &       &        何11:1 \\
    \midrule
    \multirow{2}{*}{耶穌受浸} & 受聖靈所膏 & 3:16 & \multicolumn{1}{p{4.39em}|}{1:10} & \multicolumn{1}{p{4.165em}|}{3:22} &        賽42:1，11:2-3，61:1 \\
\cmidrule{2-6}    \multicolumn{1}{|r|}{} & 神的兒子  & 3:17 & \multicolumn{1}{p{4.39em}|}{1:11} & \multicolumn{1}{p{4.165em}|}{3:22} &        詩2:7 \\
    \midrule
        耶穌被拿撒勒人棄絕 & 報好消息  &   & \multicolumn{1}{r|}{} &   4:18-21            & 賽61:1 \\
    \midrule
    耶穌又被拿撒勒人厭棄 & 被人藐视厭棄 &  13:54-58 &  6:1-6 & \multicolumn{1}{p{4.165em}|}{ 4:22-30}        & 赛53:3 \\
    \midrule
    耶穌在迦百農傳道 & 加利利地的大光 &  4:13-17 & \multicolumn{1}{r|}{} &              & 賽9:1-2 \\
    \midrule
    耶穌醫治各樣的病人 & 代替軟弱擔當疾病 &  8:16-17 & \multicolumn{1}{r|}{} &       &         賽53:4 \\
    \midrule
  %耶穌醫治長大麻風的人 & 大麻風得潔淨 &  8:2-4 &  1:40-45 & \multicolumn{1}{p{4.165em}|}{ 5:12-16} &       & 賽35:5-6 \\   
    耶穌退到海邊去 & 不傳名   &  12:18-20 & \multicolumn{1}{r|}{} &       &        賽42:1-3 \\
   
    \bottomrule
    \end{tabular}%
      \caption{舊約聖經預言的應驗（1）}
  \label{tab:addlabel}%
\end{table}%

\begin{table}[htbp]
  \centering
    \begin{tabular}{|p{11.72em}|p{7.11em}|p{4.445em}|p{5.39em}|l|l|p{4.39em}|}
    \toprule
        \multicolumn{1}{|l|}{\textbf{事件}} & \multicolumn{1}{l|}{\textbf{內容}} & \multicolumn{1}{l|}{\textbf{馬太福音}} & \multicolumn{1}{l|}{\textbf{馬可福音}} & \multicolumn{1}{l|}{\textbf{路加福音}} & \multicolumn{1}{l|}{\textbf{約翰福音}} &   \multicolumn{1}{l|}{\textbf{舊約聖經}} \\
    \midrule
    耶穌講天國的比喻 & 開口說比喻 &  13:1-53 &  4:1-34 &    &    & 詩78:2 \\
      \midrule
    耶穌潔淨聖殿 & 為父神的殿心焦如火燒 &  & \multicolumn{1}{r|}{} &       &   2:17     & 詩69:9 \\
    \midrule
    耶穌與尼哥底母談道 & 賜人永生  &   & \multicolumn{1}{r|}{} &       &     3:15-16   & 詩89:48 \\
    \midrule
    耶穌與婦人傳道 & 基督必將一切的事告訴人 &  & \multicolumn{1}{r|}{} &       &    4:25    & 申18:18 \\
 
    \midrule
    耶穌醫治兩個瞎子 & \multirow{4}[8]{*}{瞎子看見} &  9:27-31 & \multicolumn{1}{r|}{} &       &       & \multirow{4}[8]{*}{賽35:5-6} \\
\cmidrule{1-1}\cmidrule{3-6}    耶穌使伯賽大瞎子看見 & \multicolumn{1}{r|}{} &       &  8:22-26 &       &       & \multicolumn{1}{c|}{} \\
\cmidrule{1-1}\cmidrule{3-6}    耶穌醫治生來瞎眼者 & \multicolumn{1}{r|}{} &       & \multicolumn{1}{r|}{} &       & \multicolumn{1}{p{4.055em}|}{ 9:1-41} & \multicolumn{1}{c|}{} \\
\cmidrule{1-1}\cmidrule{3-6}    耶穌醫治耶利哥的瞎子 & \multicolumn{1}{r|}{} &  20:29-34 &  10:46-52 & \multicolumn{1}{p{4.165em}|}{ 18:35-43} & \multicolumn{1}{p{4.055em}|}{ } & \multicolumn{1}{c|}{} \\
    \midrule
    耶穌醫聾啞其他病人 &  耳聾舌結的人 & 15:29-31 &  7:31-37 &       &       & 賽35:5-6 \\
    \midrule
    耶穌登山變像 & 神的兒子  &  17:5 &  9:7  & \multicolumn{1}{p{4.165em}|}{ 9:35} &       & \multicolumn{1}{l|}{詩2:2} \\
    \midrule
    門徒爭論誰為大 & 贖罪祭   &  20:28 &  10:45 &       &       & 賽53:10 \\ \midrule
        \multirow{2}{*}{耶稣騎驢進城} & 騎驢進耶路撒冷 &  21:1-7 &  11:1-7 &  19:29-35 & \multicolumn{1}{p{4.055em}|}{ 12:14-16} & 亞9:9 \\
\cmidrule{2-7}    \multicolumn{1}{|r|}{} & 榮耀進城  &  21:8-10 &  11:8-10 &  19:36-38 & \multicolumn{1}{p{4.055em}|}{ 12:12-13} & 詩24:7-10 \\
    \midrule
 
耶穌第二次潔淨聖殿    & 小孩的讚美 &  21:15-16 &   &   &       & 詩8:2 \\
    \midrule
    耶穌論他的權柄 & 匠人所棄房角石 & 21:42 & 12:10 & 20:17 &       & 詩118:22 \\
    \midrule
    耶穌論基督與大衛的關係 & 基督是主  & 22:44 & \multicolumn{1}{r|}{} & 20:42 &       & 詩110:1 \\
        \midrule  
     耶穌最後公開的講論 & 法利賽人不信耶穌 &  & \multicolumn{1}{r|}{} & \multicolumn{1}{r|}{} & \multicolumn{1}{p{4.055em}|}{12:37-39} & 賽6:10，53:1 \\
    \midrule
    猶大與祭司長商量賣主 & 三十塊錢  & 26:15 & \multicolumn{1}{r|}{} & \multicolumn{1}{r|}{} &       & 亞11:12 \\
    \midrule
    耶穌為門徒洗腳 & 被一同吃飯的人出賣 &       & \multicolumn{1}{r|}{} & \multicolumn{1}{r|}{} & \multicolumn{1}{p{4.055em}|}{13:18} & 詩41:9 \\
    \midrule
    耶穌最後的談話 & 世人無故恨主與門徒 &  & \multicolumn{1}{r|}{} & \multicolumn{1}{r|}{} & \multicolumn{1}{p{4.055em}|}{15:25} & 詩69:4 \\
    \midrule
    耶穌被捉拿 & 擊打牧人羊就分散 & 26:31，54-56 & 14:27，49-50 & \multicolumn{1}{r|}{} &       & 亞13:7 \\
    \midrule
    \multirow{2}[4]{*}{猶大自盡} & 年日短少別人得他的職分 & 27:5 & \multicolumn{1}{r|}{} & \multicolumn{1}{r|}{} & \multicolumn{1}{p{4.055em}|}{17:12} & 詩109:9 \\
\cmidrule{2-7}    \multicolumn{1}{|r|}{} & 三十塊錢給窯戶 & 27:7 & \multicolumn{1}{r|}{} & \multicolumn{1}{r|}{} &       & 亞11:13 \\
    \bottomrule
    \end{tabular}%
      \caption{舊約聖經的預言應驗（2）}
  \label{tab:addlabel}%
\end{table}%

 
 
 \begin{table}[htbp]
  \centering

    \begin{tabular}{|p{8.72em}|p{7.11em}|p{4.445em}|p{4.39em}|p{4.165em}|r|p{8.39em}|}
    \toprule
    \multicolumn{1}{|l|}{\textbf{事件}} & \multicolumn{1}{l|}{\textbf{內容}} & \multicolumn{1}{l|}{\textbf{馬太福音}} & \multicolumn{1}{l|}{\textbf{馬可福音}} & \multicolumn{1}{l|}{\textbf{路加福音}} & \multicolumn{1}{l|}{\textbf{約翰福音}} &   \multicolumn{1}{l|}{\textbf{舊約聖經}} \\
    \midrule
    \multirow{3}{*}{耶穌被抓受審} & 未行強暴口無詭詐 & 26:60，\newline{27：18-19}23-24， & 14:55 & 23:4，14-15，23：22 & \multicolumn{1}{p{4.055em}|}{18:38，\newline{19：6}} & 賽53:9 \\
\cmidrule{2-7}    \multicolumn{1}{|c|}{} & 受審不言  & 26:63，27：12-14 & 14：61，15：5 & 23:9 & \multicolumn{1}{p{4.055em}|}{19:9} & 賽53:7 \\
\cmidrule{2-7}    \multicolumn{1}{|c|}{} & 被人辱罵欺凌羞辱 & 26:67-68，27：39-40 & 14:65，15：16-20，29 & 22:63-65 & \multicolumn{1}{p{4.055em}|}{19:1-3} & 詩69:19 \\

    \midrule
    耶穌被交給人釘死 & 刑罰鞭打  & 26:67，27：26 & 15:15 & \multicolumn{1}{r|}{} & \multicolumn{1}{p{4.055em}|}{19:1} & 賽53:5，53:8 \\
    \midrule
    \multirow{2}[4]{*}{耶穌被釘十字架} & 手腳被扎  &  & \multicolumn{1}{r|}{} & 23:33 & \multicolumn{1}{p{4.055em}|}{20:25} & 詩22:16 \\
\cmidrule{2-7}    \multicolumn{1}{|r|}{} & 列在罪犯中 & 27:38 & 15:27 & 23:33 & \multicolumn{1}{p{4.055em}|}{19:18} & 賽53:12 \\
    \midrule
    耶穌為人代禱 & 為罪犯代求 &  & \multicolumn{1}{r|}{} & 23:34 &       & 賽53:12 \\
    \midrule
  %  耶穌被眾人譏笑 & 被人辱罵欺凌羞辱 & 27:39-40 & 15:29 & \multicolumn{1}{r|}{} &       & 詩69:19 \\
    %\midrule
    兵丁分衣  & 分外衣裏衣拈鬮 &       & \multicolumn{1}{r|}{} & 23:24 & \multicolumn{1}{p{4.055em}|}{19:23-24} & 詩22:18 \\
    \midrule
    \multirow{6}[12]{*}{耶穌斷氣前的光景} & 被神棄絕  & 27:46 & \multicolumn{1}{r|}{} & \multicolumn{1}{r|}{} &       & 詩22:1，89：38 \\
\cmidrule{2-7}    \multicolumn{1}{|r|}{} & 苦膽調酒  & 27:34 & \multicolumn{1}{r|}{} & \multicolumn{1}{r|}{} &       & 詩69:21 \\
\cmidrule{2-7}    \multicolumn{1}{|r|}{} & 口渴喝醋  & 27:48 & \multicolumn{1}{r|}{} & 23:36 &       & 詩69:21 \\
\cmidrule{2-7}    \multicolumn{1}{|r|}{} & 交靈魂給神 &  & \multicolumn{1}{r|}{} & 23:46 & \multicolumn{1}{p{4.055em}|}{19:30} & 詩31:5 \\
\cmidrule{2-7}    \multicolumn{1}{|r|}{} & 將命傾倒  & 27:50 & 15:37 & 23:46 &       & 賽53:12 \\
\cmidrule{2-7}    \multicolumn{1}{|r|}{} & 青年離世  &       & \multicolumn{1}{r|}{} & \multicolumn{1}{r|}{} &       & 詩89:45 \\
    \midrule
    \multirow{2}[4]{*}{耶穌肋旁被扎} & 骨頭不折斷 &  & \multicolumn{1}{r|}{} & \multicolumn{1}{r|}{} & \multicolumn{1}{p{4.055em}|}{19:31-33，36} & 出12:46，民9:12，詩34:20 \\
\cmidrule{2-7}    \multicolumn{1}{|r|}{} & 救主被扎  & & \multicolumn{1}{r|}{} & \multicolumn{1}{r|}{} & \multicolumn{1}{p{4.055em}|}{ 19:34，37} & 亞12:10 \\
    \midrule
    耶穌被安葬 & 與財主同葬 &  27:57-61 &  15:42-47 &  23:50-56 & \multicolumn{1}{p{4.055em}|}{ 19:38-42} & 賽53:9 \\
    \midrule
    使者宣告主已復活 & 復活    &  28:5-7 &  16:5-7 &  24:3-7 &       & 詩16:10，117：17 \\
    \midrule
    耶穌被接升天 & 上升為至高 &  & 16:19 &  24:51 &       & 賽52:13 \\
    \bottomrule  
 
      \end{tabular}%
         \caption{舊約聖經的預言應驗（3）}
  \label{tab:addlabel}%
\end{table}%

